\chapter*{Informazioni sul corso e sull'esame}
\addcontentsline{toc}{chapter}{Informazioni sul corso e sull'esame}
\markboth{Informazioni}{Informazioni}

\textbf{Docente:} Marta Pascoal (DEIB, edificio 20, II piano), \texttt{marta.brazpascoal@polimi.it}, ricevimento su appuntamento.

\textbf{Orari:} martedì 14:15--16:15 (aula 3.0.2, lezione); venerdì 10:15--12:15 (aula 25.S.2, lezione/esercitazione); venerdì 13:15--15:15 (aula 3.0.2, laboratorio/recuperi). Tutto registrato e in streaming.

\textbf{Esame.}
\begin{itemize}[nosep]
  \item Scritto (in presenza), max 30 punti: esercizi, \emph{domande teoriche} e conoscenza del \emph{linguaggio di modellazione} visto in laboratorio.
  \item Niente libri, appunti, calcolatrici o dispositivi (possederli comporta penalità o annullamento).
  \item Punteggio $\le 10$ $\Rightarrow$ respinto. Nessuna prova in itinere. Eventuale orale a discrezione del docente, \emph{non integrativo} (può anche abbassare il voto).
  \item \textbf{Quesiti brevi}: 4 quiz da 15 minuti, all'inizio della lezione del martedì, solo in presenza; valgono in tutto 3 punti (si scarta il peggiore); validi solo per l'A.A.\ in corso.
\end{itemize}

\textbf{Quesiti brevi: come funzionano} (regole pubblicate a ottobre).
\begin{itemize}[nosep]
  \item Si svolgono su Moodle \textbf{dallo smartphone} (link su WeBeep all'inizio della prova): aggiorna prima il sistema operativo. Durata 15 minuti; niente materiale, niente calcolatrice, niente domande al docente.
  \item Formato \textbf{Vero/Falso}: ogni domanda ha più affermazioni (es.\ Q1.1--Q1.4). Le stesse risposte (``V''/``F'') vanno ricopiate sul foglio cartaceo con nome, cognome e codice persona.
  \item Ogni quesito vale tra 0 e 1 punto, frazionario in base alle risposte corrette, senza arrotondamento; si scarta il peggiore dei 4. I punti si sommano al punteggio dello scritto \emph{solo se lo scritto è sufficiente}; si arrotonda solo il totale.
  \item Per allenarti: Cap.~\ref{cap:quiz} (affermazioni V/F con soluzioni).
\end{itemize}

\textbf{Date:} quesiti brevi il 6/10, 3/11, 24/11, 15/12 (ore 14:15). Laboratori il 23/10, 6/11, 27/11, 11/12 (ven.\ 13:15). Lezioni sospese il 29/9 e il 2/10. \textbf{Primo quesito breve: martedì 6/10} (presumibilmente sugli argomenti dei Capitoli 1--3).

\textbf{Programma del corso}
\begin{enumerate}[nosep]
  \item Modellazione in Programmazione Lineare (PL): tipi di variabili, vincoli, funzioni obiettivo. \hfill \emph{(questi appunti)}
  \item Cenni di complessità computazionale. \hfill \emph{(questi appunti)}
  \item Proprietà della PL e metodo del simplesso.
  \item Metodi per la Programmazione Lineare Intera (PLI).
  \item Problemi di ottimizzazione su grafo.
  \item Laboratorio: linguaggio di modellazione e solutori.
\end{enumerate}

\textbf{Mappa delle lezioni coperte}

\begin{tabular}{@{}llp{0.78\textwidth}@{}}
\toprule
Lezione & Data & Contenuto \\ \midrule
1 & 15/9 & Introduzione, esempio delle sale operatorie (\S\ref{sec:sale}) \\
2 & 18/9 & Soluzioni ammissibili/ottime, forma parametrica e compatta; coltivatore, dieta (\S\ref{sec:def}, \S\ref{sec:colt}, \S\ref{sec:dieta}) \\
3 & 22/9 & Investimento, zaino, big-$M$, vincoli logici, blending, localizzazione bottleneck (\S\ref{sec:zaino}--\S\ref{sec:bottleneck}) \\
4 & 25/9 & Variabili libere e discrete (\S\ref{sec:libere}, \S\ref{sec:discrete}); complessità (Cap.~\ref{cap:compl}) \\
5 & 25/9 & Esercitazione 1: modelli, in aula esercizi 1, 6, 8, 9 (documento \emph{Esercitazioni FRO}) \\
6+ & dal 6/10 & Programmazione Lineare: forme, geometria, teorema fondamentale (Cap.~\ref{cap:pl}) \\
\bottomrule
\end{tabular}
